\documentclass[11pt]{article}
\usepackage[margin=1in]{geometry}
\usepackage{amsmath,amsthm,amssymb}
\usepackage[hidelinks]{hyperref}
\newtheorem*{theorem}{Theorem}
\title{Primes Between Consecutive Cubes\\[3pt]\large A consequence of the Riemann hypothesis}
\author{}
\date{22 November 2026}
\begin{document}
\maketitle
\begin{abstract}
We give a short derivation of the classical result that every open interval
between consecutive positive cubes contains a prime. We apply Dudek's
explicit short-interval theorem to the upper cube and check the 62
remaining cases. The accompanying Lean file verifies this reduction from
the explicit interval theorem, not yet the entire chain from the
recent Lean proof of the Riemann hypothesis.
\end{abstract}

\begin{theorem}
For every integer $n\ge1$, there is a prime $p$ with
$n^3<p<(n+1)^3$.
\end{theorem}

This is a known consequence of the Riemann hypothesis (RH), proved by
Caldwell and Cheng \cite{CC}. Here we give a short reduction using a later
bound of Dudek. The mathematical argument below uses classical RH as
established in the formal development announced by Alp\"oge \cite{A};
the Lean and Mathlib communities independently verified the statement
and proof \cite{V}. No generalized RH is needed.

\begin{proof}
Dudek \cite[Theorem~2]{D} shows under RH that for each real $x\ge2$ there
is a prime in
\[
\left(x-\frac4\pi\sqrt{x}\log x,\ x\right].
\]
For $n\ge63$, set $m=n+1$ and $s=\sqrt m\ge8$.
The inequality $\log t\le t-1$ for $t>0$ gives
\[
\log m=2\bigl(3\log 2+\log(s/8)\bigr)
 \le \frac{11}{5}+\frac{s}{4}\le\frac{3s}{5},
\]
using $\log 2<7/10$. Hence at $x=m^3$ the interval length satisfies
\[
\frac4\pi\sqrt{m^3}\log(m^3)
=\frac{12}{\pi}\,ms\log m
< \frac{12}{5}m^2
< 3m^2-3m+1=m^3-(m-1)^3 .
\]
The first strict inequality uses $\pi>3$ and $m>1$; the second follows
from $m\ge64$. Thus Dudek's prime is strictly above $(m-1)^3$. It
cannot equal $m^3$, which is composite.

For $1\le n\le62$, the Lean source supplies a prime witness for each
interval and checks its primality and the two strict inequalities in the
kernel.
\end{proof}

\paragraph{Formal verification status.}
\texttt{CubePrime.from\_short\_interval} proves the reduction and all
finite checks with \texttt{ShortIntervalBound} as an explicit premise.
Hubbard's formalization of Dudek's estimate \cite{H} takes, in addition
to RH, a zero-counting premise \texttt{HSWN70}. The cited zero-count
estimate is a mathematical theorem \cite{HSW}; its exact premise has
not yet been discharged in this Lean composition. The two developments
also use different Lean versions. The source therefore does not yet
prove the unconditional theorem end to end in Lean.

\begin{thebibliography}{9}\small
\bibitem{CC} C. K. Caldwell and Y. Cheng, \emph{Determining Mills'
constant and a note on Honaker's problem}, J. Integer Seq. 8 (2005),
Article 05.4.1, Lemma 5.
\href{https://cs.uwaterloo.ca/journals/JIS/VOL8/Caldwell/caldwell78.html}{Journal}.
\bibitem{D} A. W. Dudek, \emph{On the Riemann hypothesis and the
difference between primes}, Int. J. Number Theory 11 (2015), 771--778.
\href{https://arxiv.org/abs/1402.6417v2}{arXiv:1402.6417v2}.
\bibitem{A} L. Alp\"oge, \emph{Announcement of a Lean proof of the
Riemann hypothesis}, 20 November 2026.
\href{https://alpo.ge/remark.pdf}{Announcement};
\href{https://github.com/leanprover-community/riemann-hypothesis/blob/58afe6f4b8167e8945f7c82954b0fc477541f9d0/src/Main.lean}{pinned Lean statement}.
\bibitem{V} Lean and Mathlib communities, \emph{Lean verification of
the Riemann hypothesis}, 21 November 2026.
\href{https://leanprover-community.github.io/blog/posts/riemann-hypothesis/}{Verification statement}.
\bibitem{H} S. Hubbard, \emph{An explicit prime gap under the Riemann
hypothesis with an elementary prime power window and an updated zero count}.
\href{https://hexagonmath.org/2611.00157v2}{Hexagon 2611.00157v2};
\href{https://github.com/sthubbar/goldbach-lean/tree/93b86c05bb441be2ac0b3f06cc4c8624f8e8d0eb/entries/dudek2014}{pinned Lean source}.
\bibitem{HSW} E. Hasanalizade, Q. Shen and P.-J. Wong,
\emph{Counting zeros of the Riemann zeta function}.
\href{https://arxiv.org/abs/2107.06506}{arXiv:2107.06506}.
\end{thebibliography}
\end{document}
